\documentclass[conference]{IEEEtran}
\usepackage{cite,amsmath,booktabs,graphicx,url}
\title{Reproducible Criterion-Level Clinical Trial Screening with an Open-Weight Transformer}
\author{\IEEEauthorblockN{Student Author}\IEEEauthorblockA{Department and institution details to be completed}}
\begin{document}
\maketitle
\begin{abstract}
This paper describes an educational prototype for classifying individual clinical-trial criteria against patient descriptions and returning sentence-level evidence. The implementation uses the required TrialGPT Criterion Annotations data, a deterministic patient-grouped split, an open-weight autoregressive model interface, parameter-efficient supervised adaptation, and explicit implementations of cached decoding, tiled online-softmax attention, and greedy speculative decoding. A Colab Tesla T4 run completed supervised adaptation and validation, and a synthetic attention check measured numerical agreement and runtime. The held-out test and decoder hardware benchmarks remain pending. The reported validation results and their limitations are summarized in this paper; no held-out or decoder performance claims are made.
\end{abstract}
\begin{IEEEkeywords}clinical trial matching, language models, evidence retrieval, KV cache, online softmax, speculative decoding\end{IEEEkeywords}

\section{Introduction}
Trial eligibility criteria are expressed as natural-language constraints, while patient records contain a mixture of relevant findings and missing information. TrialGPT studies criterion-level matching and trial ranking with language models \cite{trialgpt}. This project implements an independent, locally executable educational system and focuses on traceable criterion labels, sentence IDs, and inference mechanics. It is not intended for clinical decisions.

\section{Related Work}
TrialGPT motivates separating criterion matching from downstream trial ranking \cite{trialgpt}. Low-Rank Adaptation (LoRA) reduces the number of trainable parameters through low-rank updates \cite{lora}. Transformer inference can reuse prior key/value states during autoregressive decoding. FlashAttention demonstrates that tiling with online softmax can compute exact attention while reducing memory traffic \cite{flashattn}. Speculative decoding drafts several tokens and verifies them with a target model while preserving greedy output under the corresponding greedy verification rule \cite{specdecode}. The selected Qwen2.5 model card documents its decoder architecture, parameter count, tokenizer and context length \cite{qwen}.

\section{Materials and Methods}
\subsection{Dataset and split}
The benchmark is \texttt{ncbi/TrialGPT-Criterion-Annotations}, containing 1,015 criterion annotations and expert labels/evidence IDs \cite{dataset}. The preprocessing allowlists patient note, criterion type/text, expert category, and expert sentence IDs, excluding generated GPT-4 fields. Unique patient identifiers are sorted, permuted with NumPy \texttt{default\_rng(42)}, then assigned in groups of 37/8/8 to train, validation, and test. The program verifies coverage and non-overlap. Test examples are not materialized during preparation.

\subsection{Model and adaptation}
The default target is Qwen2.5-1.5B-Instruct: a 1.54B decoder-only model with RoPE and grouped-query attention, advertised 32,768-token context and Apache-2.0 license \cite{qwen}. Qwen2.5-0.5B-Instruct (0.49B) is the draft candidate; the Qwen2.5 family uses a unified tokenizer, and code checks actual vocabularies before speculative decoding. Phi-3.5-mini-instruct (3.8B, MIT, 128K context) is a larger alternative with higher expected weight and optimizer memory requirements \cite{phi}. The repository environment is CPU-only (PyTorch 2.14.1), while the reported fine-tuning run was performed on a Colab Tesla T4. The training script applies LoRA to attention projections with rank 16, alpha 32 and dropout 0.05. It trains for three epochs with learning rate $2\times10^{-4}$, batch size one, eight gradient-accumulation steps and 4096 token maximum. Loss is masked on prompt tokens and applied to the serialized answer; checkpoint selection uses validation loss. Colab required FP16 on the T4 and a Transformers 5 compatibility edit; these edits were made only in the runtime and must be reconciled for exact repository reproduction.

The model receives the patient note, criterion type and criterion text. It must return one of six exact expert categories: included, not included, not enough information, excluded, not excluded, or not applicable, plus sentence IDs and a concise explanation. Only IDs present in the parsed source note are retained. Truncation preserves the suffix of the prompt under the token budget; truncation can discard earlier note content and therefore is a limitation to quantify.

\subsection{Evaluation and inference algorithms}
Classification uses accuracy, macro-F1, per-category precision/recall/F1 and confusion matrices overall and by inclusion/exclusion. Evidence is scored by sentence-ID micro and macro precision/recall/F1; empty/empty per-record evidence sets receive one for precision and recall, while one-sided empty sets receive zero. A patient-level outcome, if shown, is an explicitly stated rule aggregation and is not a dataset label.

The cached decoder prefills the prompt once, then forwards the newest token with the returned cache, an absolute position ID, and an attention mask covering cached and current positions. The uncached reference recomputes the full sequence each step. Tiled attention keeps running row maxima $m$, softmax normalizers $l$, and weighted-value accumulators $o$ while processing key blocks; outputs are checked against a direct scaled dot-product implementation. Greedy speculative decoding drafts a token block, verifies candidates with target logits, accepts the matching prefix, appends the target correction at first disagreement, and synchronizes cache state by re-prefilling emitted context. This conservative re-prefill favors clear state correctness and can reduce speed.

\subsection{Application}
The Streamlit application accepts numbered patient text and free-text inclusion/exclusion criteria, displays criterion category, explanation and evidence text, exposes a transparent aggregation rule and downloads a JSON report. It includes an educational-use notice and requires the local model to load successfully.

\section{Experimental Setup and Results}
Training, validation evaluation, and a synthetic attention benchmark have been run in Colab; the final test split remains untouched. The same 206 validation annotations were evaluated with the base model and adapter. Table I reports observed classification and evidence metrics. The adapter increased accuracy from 0.262 to 0.723 and macro-F1 from 0.111 to 0.356. The increase in accuracy is not a complete picture: the adapter predicted \emph{not excluded} for all 140 exclusion examples, giving perfect recall for that majority class but failing to recover the two gold \emph{excluded} and eleven \emph{not applicable} cases. Its macro-F1 on exclusion criteria is only 0.140. This class bias motivates error review and careful interpretation; test data must not be used to tune it. Evidence micro-F1 increased from 0.000 to 0.429. Row-level evaluation artifacts remain in the Colab runtime and are not committed because they contain patient notes. On six small synthetic causal/noncausal attention cases, tiled outputs matched the reference with maximum absolute error $5.36\\times10^{-7}$, but the tiled Python implementation was 2.58--51.98 times slower; these measurements do not establish an advantage on model workloads. KV-cache and speculative decoder benchmarks remain pending; they will compare output token IDs, synchronized latency, throughput, peak accelerator memory, and draft acceptance.

\begin{table}[ht]
\caption{Validation results; test and inference benchmarks pending.}
\centering
\begin{tabular}{lrr}\toprule Measure & Base & LoRA adapter\\\midrule Accuracy & 0.262 & 0.723\\Macro-F1 & 0.111 & 0.356\\Evidence micro-F1 & 0.000 & 0.429\\Inclusion accuracy & 0.424 & 0.727\\Inclusion macro-F1 & 0.116 & 0.250\\Exclusion accuracy & 0.186 & 0.721\\Exclusion macro-F1 & 0.080 & 0.140\\\bottomrule\end{tabular}
\end{table}

\section{Limitations and Failure Analysis}
The task's 53-patient sample is small; validation shows a strong class bias despite improved aggregate accuracy. Criterion labels are not independent across patients. In a targeted review of eight criterion-level errors drawn from only three patient notes, all eight gold `not enough information` exclusion labels were predicted as `not excluded`. Several predictions cited unrelated sentences; another missed the annotated uterine-imaging evidence and cited symptoms instead. This small, clustered sample illustrates failure modes but cannot estimate their prevalence. Free-form generation can produce malformed or unsupported outputs despite JSON prompting and ID validation. Evidence annotations can be empty, and absent evidence must not be interpreted as proof of eligibility or ineligibility. The held-out test, inference benchmarks, report PDF, screenshots, model artifact publication, and demo recording remain pending.

\section{Conclusion}
The repository provides a reproducible implementation scaffold covering data integrity, supervised adaptation, evidence-aware evaluation, model inference algorithms, and a browser application. Initial validation results show gains over the base model alongside a substantial class imbalance failure; broader conclusions await patient-diverse error review, decoder hardware benchmarks, and the one-time held-out test. The system remains an educational prototype and cannot determine enrollment eligibility.

\section*{AI Use Acknowledgment}
AI assistance supported repository scaffolding and implementation. The student must review and understand the code, perform the experiments, verify sources, and disclose AI use in the submitted work.

\begin{thebibliography}{00}
\bibitem{trialgpt} Q. Jin et al., ``Matching patients to clinical trials with large language models,'' \emph{Nature Communications}, vol. 15, 2024, doi: 10.1038/s41467-024-53081-z.
\bibitem{dataset} NCBI, ``TrialGPT Criterion Annotations,'' Hugging Face Datasets. [Online]. Available: \url{https://huggingface.co/datasets/ncbi/TrialGPT-Criterion-Annotations}. Accessed: Nov. 10, 2026.
\bibitem{qwen} Qwen Team, ``Qwen2.5-0.5B-Instruct Model Card.'' [Online]. Available: \url{https://huggingface.co/Qwen/Qwen2.5-0.5B-Instruct}. Accessed: Nov. 10, 2026.
\bibitem{phi} Microsoft, ``Phi-3.5-mini-instruct Model Card.'' [Online]. Available: \url{https://huggingface.co/microsoft/Phi-3.5-mini-instruct}. Accessed: Nov. 10, 2026.
\bibitem{lora} E. J. Hu et al., ``LoRA: Low-Rank Adaptation of Large Language Models,'' in \emph{Proc. ICLR}, 2022.
\bibitem{flashattn} T. Dao et al., ``FlashAttention: Fast and Memory-Efficient Exact Attention with IO-Awareness,'' in \emph{Advances in Neural Information Processing Systems}, 2022.
\bibitem{specdecode} Y. Leviathan, M. Kalman, and Y. Matias, ``Fast Inference from Transformers via Speculative Decoding,'' in \emph{Proc. ICML}, PMLR, vol. 202, 2023.
\end{thebibliography}
\end{document}
